
\setcounter{footnote}{64}

\section*{\SGAInlineRefTarget{sga3:iii:section:3}3. Infinitesimal extensions of a group scheme}
\addcontentsline{toc}{section}{3. Infinitesimal extensions of a group scheme}

Still using the notation of \S0 (\(S,\mathscr I,\mathscr J\), and so on),
let \(X\) be an \(S\)-scheme, and suppose that \(X_{\mathscr J}\) is
endowed with a group structure. We propose to determine the \(S\)-group
structures on \(X\) that induce the given structure on
\(X_{\mathscr J}\).

From now on, assume that \(X\) is flat over \(S\). Let \(\mathscr C\) be
the category of \(S\)-schemes flat over \(S\); thus
\(X\in\operatorname{Ob}\mathscr C\). Denote by \(Y\), respectively \(M\),
the functor on \(\mathscr C\) defined by \(X^+\), respectively \(L'_X\).
The canonical morphism \(p_X:X\to X^+\) defines a morphism in
\(\widehat{\mathscr C}\)
\[
 p:X\longrightarrow Y,
\]
and the action of \(L'_X\) on \(X\) in
\(\widehat{(\mathrm{Sch})/S}\) defines an action of \(M\) on \(X\) in
\(\widehat{\mathscr C}\). One checks at once that \(X\) is thereby
formally principal homogeneous under \(M_Y\) over \(Y\) (cf. \SGARefLink{sga3:iii:proposition:0-2:item:i}{0.2(i)}
and \SGARefLink{sga3:iii:remark:0-4}{0.4}).

The action of \(X^+\) on \(L'_X\) defined in \SGARefLink{sga3:iii:proposition:0-8}{0.8} (denoted
\(\operatorname{Ad}\) there) defines an action, denoted \(\widetilde f\),
of \(Y\) on \(M\). On the other hand, by \SGARefLink{sga3:iii:remark:0-5}{0.5},
\[
 \operatorname{Hom}_{\widehat{\mathscr C}}(Z,M)
 \simeq \operatorname{Hom}_{S_0}(Z_0,L_0),
 \qquad Z\in\operatorname{Ob}\mathscr C,
\]
where \(L_0\) is the functor defined in \SGARefLink{sga3:iii:remark:0-5}{0.5}.

\medskip
\noindent\SGAInlineRefTarget{sga3:iii:lemma:3-1}\textbf{Lemma 3.1.}
\begin{enumerate}
\item[(i)] \textit{Condition \((+)_{n}\) of \SGARefLink{sga3:iii:subsection:1-3}{1.3} holds for every positive
integer \(n\).}
\item[\SGAInlineRefTarget{sga3:iii:lemma:3-1:item:ii}(ii)] \textit{Let the \(S_0\)-group \(X_0\) act on the
\(S_0\)-functor \(L_0\) through its adjoint representation. Then there is
a canonical isomorphism}
\[
 H^*(X_0,L_0)\simeq H^*(Y,M),
\]
\textit{where the first cohomology is computed in
\((\mathrm{Sch})/S_0\), and the second in \(\mathscr C\).}
\end{enumerate}

Both parts of the lemma follow from the relation
\begin{align*}
 \operatorname{Hom}_{\widehat{\mathscr C}}(Y,M)
 &\simeq
 \operatorname{Hom}_{\widehat{(\mathrm{Sch})/S_0}}
       (X^+\times_S S_0,L_0)\\
 &\simeq \operatorname{Hom}_{S_0}(X_0,L_0)\\
 &\simeq \operatorname{Hom}_{\widehat{\mathscr C}}(X,M),
\end{align*}
which follows at once from the definition of \(M\) as
\(\prod_{S_0/S}L_0\). More generally, the same relation holds with
\(X,Y\) replaced by \(X^n,Y^n\). Consequently, every morphism
\(X^n\to M\) factors uniquely through \(Y^n\), which gives \((+)_{n}\).
It also gives the relation
\(C^*(Y,M)=C^*(X_0,L_0)\), and hence \SGARefLink{sga3:iii:lemma:3-1:item:ii}{(ii)}.

We may therefore apply the constructions of \SGARefLink{sga3:iii:subsection:1-3}{1.3}. In particular:

\medskip
\noindent\SGAInlineRefTarget{sga3:iii:lemma:3-2}\textbf{Lemma 3.2.}
\textit{Let \(P:X\times_S X\to X\) be a morphism. Then \(P\) induces
the group law of \(X_{\mathscr J}\) if and only if \(P\) is an
admissible composition law on \(X\) (cf. \SGARefLink{sga3:iii:definition:1-3-2}{1.3.2}).}

Indeed, \(P\) induces the group law of \(X_{\mathscr J}\) if and only if
it lifts the group law of \(X^+\), or equivalently that of \(Y\). Thus it
remains only to show that every morphism \(P\) lifting the group law of
\(X_{\mathscr J}\) satisfies identity ((++)) of \SGARefLink{sga3:iii:definition:1-3-2:item:ii}{1.3.2(ii)}, which is
exactly what was proved in \SGARefLink{sga3:iii:proposition:0-8}{0.8}.

\medskip
\noindent\SGAInlineRefTarget{sga3:iii:proposition:3-3}\textbf{Proposition 3.3.}
\textit{Let \(S\) be a scheme and \(S_0\) a closed subscheme defined by
a nilpotent ideal. Let \(X\) be an \(S\)-scheme that is flat and either
quasi-compact or locally of finite presentation over \(S\). Let
\(P:X\times_S X\to X\) be a composition law on \(X\). Then \(P\) is a
group law if and only if the following two conditions hold:}
\begin{enumerate}
\item[\SGAInlineRefTarget{sga3:iii:proposition:3-3:item:i}(i)] \textit{\(P\) is associative;}
\item[\SGAInlineRefTarget{sga3:iii:proposition:3-3:item:ii}(ii)] \textit{\(P\) induces a group law on
\(X_0=X\times_S S_0\).}
\end{enumerate}

These conditions are clearly necessary. We prove that they are
sufficient. Suppose first that \(X\to S\) has a section. Since
\(X(S')\) is then nonempty for every \(S'\to S\), it suffices
\footnote{Editor's note: The original has been corrected by deleting
the inappropriate reference to an exercise of Bourbaki on semigroups
(cf. [BAlg], \S I.2, Exercises 9--13) and by indicating the role of the
left and right translations; see the following editor's note.}
to show that, for every \(x\in X(S')\), left and right translation by
\(x\) are isomorphisms of \(X_{S'}\).
\footnote{Editor's note: Let \(E\) be a nonempty set endowed with an
associative composition law such that every left translation
\(\ell_x\) is bijective, and fix \(x_0\in E\). There is a unique
\(e\in E\) such that \(x_0e=x_0\); then
\(x_0ex=x_0x\) implies \(ex=x\) for every \(x\in E\). Moreover, for
every \(x\) there is a unique \(x'\) such that \(xx'=e\). Suppose in
addition that there is a \(b\in E\) for which right translation
\(r_b\) is injective. Then \(xeb=xb\) gives \(xe=x\), so \(e\) is a
unit, and \(xx'x=x=xe\) gives \(x'x=e\); hence \(x'\) is both a left
and a right inverse of \(x\), and \(E\) is a group.

The hypothesis that \(r_b\) be injective is necessary: on any set
\(E\), define \(xy=y\) for all \(x,y\in E\). Every left translation is
then the identity (and the law is associative), whereas
\(r_y(E)=\{y\}\) for every \(y\); thus \(E\) is not a group when
\(|E|>1\).}

We may plainly assume \(S'=S\). The translation in question, \(t\),
induces on \(X_0\) a translation \(t_0\), which is an automorphism since
\(X_0\) is a group. The conclusion follows by flatness (SGA 1 III,
4.2).
\footnote{Editor's note: Since \(X\) and \(X_0\) have the same
underlying topological space and \(t_0\) is an automorphism, \(t\) is a
homeomorphism, hence an affine morphism; cf. Expos\'e VI\textsubscript{B},
2.9.1, or EGA IV\textsubscript{4}, 18.12.7.1. It therefore suffices to
show the following. Let \(\mathscr J\) be a nilpotent ideal of a ring
\(\Lambda\), and let \(\varphi:A\to B\) be a morphism of
\(\Lambda\)-algebras, with \(B\) flat over \(\Lambda\), such that
\(\varphi\otimes_\Lambda(\Lambda/\mathscr J)\) is bijective. Then
\(\varphi\) is bijective. By the ``nilpotent Nakayama lemma,''
\(\varphi\) is surjective; moreover, since \(B\) is flat over
\(\Lambda\), one also has
\(\operatorname{Ker}(\varphi)\otimes_\Lambda
(\Lambda/\mathscr J)=0\), whence
\(\operatorname{Ker}(\varphi)=0\). Thus \(\varphi\) is bijective.}

Now drop the assumption that \(X\) has a section over \(S\), and
suppose that there exists \(S'\to S\) such that \(X_{S'}\) has a
section over \(S'\). By what has just been proved, \(X_{S'}\) is then an
\(S'\)-group; let \(e'\) be its unit section. The inverse image of
\(e'\) under
\(\operatorname{pr}_i:S'\times_S S'\to S'\), \(i=1,2\), is the unit
section of \(X_{S''}\) for the group law obtained by pulling
\(P_{S'}\) back along \(\operatorname{pr}_i\). But since \(P\) is
``defined over \(S\),'' these two group laws coincide, and so do their
unit sections. Hence
\[
 \operatorname{pr}_1^*(e')=\operatorname{pr}_2^*(e').
\]

If \(S'\to S\) is a descent morphism (cf. Expos\'e IV, no.~2), there is
a section of \(X\) whose base extension is \(e'\), and the proof is
complete. Since \(X_X\) has a section over \(X\), namely the diagonal
section, it is now enough to prove that \(X\to S\) is a descent
morphism. It is flat and surjective, and is either quasi-compact or
locally of finite presentation; it is therefore an (fpqc) covering,
hence a descent morphism (Expos\'e IV, no.~6).

\medskip
\noindent\textbf{Remark.}
In fact, the hypothesis that \(X\to S\) be quasi-compact or locally of
finite presentation is superfluous, by the following result, which the
reader may prove as an exercise on Expos\'e IV:

Under the stated hypotheses on \(S\) and \(S_0\), if \(X\to S\) is
flat and \(X_0\to S_0\) is an (fpqc) covering, then \(X\to S\) is a
descent morphism.

\medskip
\noindent\SGAInlineRefTarget{sga3:iii:lemma:3-4}\textbf{Lemma 3.4.}
\textit{Two admissible composition laws on \(X\) are equivalent
(cf. \SGARefLink{sga3:iii:definition:1-3-5}{1.3.5}) if and only if one is obtained from the other by an
automorphism of \(X\) over \(S\) that induces the identity on
\(X_{\mathscr J}\).}

Indeed, the morphisms constructed in \SGARefLink{sga3:iii:lemma:1-3-1}{1.3.1} are precisely those in the
preceding statement, by \SGARefLink{sga3:iii:corollary:0-7}{0.7}.
\footnote{Editor's note: Indeed, by the proof of \SGARefLink{sga3:iii:corollary:0-7}{0.7}, the
\(S\)-endomorphisms of \(X\) inducing the identity on
\(X_{\mathscr J}\) are the automorphisms \(m\cdot\operatorname{id}_X\),
where \(m\) ranges over
\(M(X)=\operatorname{Hom}_S(X,M)\). For every \(S'\to S\) and every
\(x\in X(S')\), one has
\((m\cdot\operatorname{id}_X)(x)=m(x)\cdot x\). By the proof of \SGARefLink{sga3:iii:lemma:3-1}{3.1},
each \(m:X\to M\) factors uniquely through a morphism
\(h:Y=X^+\to M\); hence \(m\cdot\operatorname{id}_X\) is the
automorphism \(u_h\) introduced in \SGARefLink{sga3:iii:lemma:1-3-1}{1.3.1}. The lemma now follows from
the definition of equivalence; cf. \SGARefLink{sga3:iii:lemma:1-3-4}{1.3.4} and \SGARefLink{sga3:iii:definition:1-3-5}{1.3.5}.}

Taking account of all the preceding results, \SGARefLink{sga3:iii:proposition:1-3-6}{Proposition 1.3.6} gives:

\medskip
\noindent\SGAInlineRefTarget{sga3:iii:theorem:3-5}\textbf{Theorem 3.5.}
\textit{Let \(S\) be a scheme, and let \(\mathscr I\) and
\(\mathscr J\) be two ideals on \(S\) such that
\(\mathscr I\supset\mathscr J\) and
\(\mathscr I\mathscr J=0\). Let \(S_0\) and \(S_{\mathscr J}\) be the
closed subschemes of \(S\) that they define. Let \(X\) be an
\(S\)-scheme flat over \(S\) (and either locally of finite presentation
or quasi-compact over \(S\)), and let \(X_0\) and
\(X_{\mathscr J}\) be the schemes obtained by base change. Suppose
\(X_{\mathscr J}\) is endowed with an \(S_{\mathscr J}\)-group
structure, and let \(L_0\) denote the \(S_0\)-functor in commutative
groups defined by}
\[
 \operatorname{Hom}_{S_0}(T,L_0)
 =\operatorname{Hom}_{\mathscr O_T}
 \left(
  \omega^1_{X_0/S_0}\otimes_{\mathscr O_{S_0}}\mathscr O_T,
  \mathscr J\otimes_{\mathscr O_{S_0}}\mathscr O_T
 \right),
\]
\textit{on which \(X_0\) acts through its adjoint representation.}
\begin{enumerate}
\item[\SGAInlineRefTarget{sga3:iii:theorem:3-5:item:i}(i)] \textit{An \(S\)-group structure exists on \(X\) that
induces the given structure on \(X_{\mathscr J}\) if and only if the
following conditions hold:}
  \begin{enumerate}
  \item[\((i_1)\)] \textit{there is a morphism of \(S\)-schemes
  \(X\times_S X\to X\) inducing the group law of
  \(X_{\mathscr J}\);}
  \item[\((i_2)\)] \textit{a certain obstruction class in
  \(H^3(X_0,L_0)\), canonically defined by \(X\) and the group law of
  \(X_{\mathscr J}\), is zero.}
  \end{enumerate}
\item[\SGAInlineRefTarget{sga3:iii:theorem:3-5:item:ii}(ii)] \textit{If the conditions in \SGARefLink{sga3:iii:theorem:3-5:item:i}{(i)} hold, the set
\(\mathcal E\) of group laws on \(X\) inducing the given law on
\(X_{\mathscr J}\) is principal homogeneous under
\(Z^2(X_0,L_0)\). Moreover, the quotient of \(\mathcal E\) by the
\(S\)-automorphisms of \(X\) inducing the identity on
\(X_{\mathscr J}\) is principal homogeneous under \(H^2(X_0,L_0)\).}
\end{enumerate}

\footnote{Editor's note: The following argument has been added.}
Indeed, by \SGARefLink{sga3:iii:lemma:3-2}{3.2} every morphism of \(S\)-schemes
\(f:X\times_S X\to X\) inducing the group law of
\(X_{\mathscr J}\) is an admissible composition law on \(X\). By
\SGARefLink{sga3:iii:proposition:1-3-6:item:i}{1.3.6(i)}, the existence of an associative admissible composition law
\(P:X\times_S X\to X\) is equivalent to the vanishing of a certain
class \(c(f)\in H^3(X_0,L_0)\); in that case, \SGARefLink{sga3:iii:proposition:3-3}{3.3} shows that \(P\) is a
group law. This proves \SGARefLink{sga3:iii:proposition:3-3:item:i}{(i)}, and \SGARefLink{sga3:iii:proposition:3-3:item:ii}{(ii)} then follows from \SGARefLink{sga3:iii:proposition:3-3}{3.3} and
\SGARefLink{sga3:iii:proposition:1-3-6:item:ii}{1.3.6(ii)}.

\medskip
\noindent\SGAInlineRefTarget{sga3:iii:remark:3-5-1}\textbf{Remark 3.5.1.}
\footnote{Editor's note: This remark, analogous to \SGARefLink{sga3:iii:remark:4-5-1}{4.5.1}, has been
added to introduce the notation \(\delta(\mu,\mu')\), or
\(\delta(X,X')\), used in \SGARefLink{sga3:iii:paragraph:4-38}{4.38}. Consequently, the part of \SGARefLink{sga3:iii:theorem:3-5:item:ii}{3.5(ii)}
concerning \(\mathcal E\) itself has also been added.}
If \(\mu\) and \(\mu'\) are group laws on \(X\) inducing the given law
on \(X_{\mathscr J}\), one obtains a cocycle
\(\delta(\mu,\mu')\in Z^2(X_0,L_0)\). The chosen sign convention is
\(\mu'=\delta(\mu,\mu')\cdot\mu\); that is, for every \(S'\to S\) and
\(x,y\in X(S')\),
\[
 \mu'(x,y)=\delta(\mu,\mu')(x_0,y_0)\cdot\mu(x,y).
\]
\footnote{Editor's note: We follow the sign conventions of the
original in order to have, in \SGARefLink{sga3:iii:paragraph:4-38:item:5}{4.38(5)}, the equality
\(\partial^1\bar d(X,X')=\bar\delta(X,X')\); see also Editor's
Note 54.}

Write \(\bar\delta(\mu,\mu')\) for the image of
\(\delta(\mu,\mu')\) in \(H^2(X_0,L_0)\). Finally, if \(X\) endowed
with \(\mu\), respectively \(\mu'\), is denoted simply by \(X\),
respectively \(X'\), write \(\delta(X,X')\) in place of
\(\delta(\mu,\mu')\), and similarly for \(\bar\delta(X,X')\).

\medskip
\noindent\SGAInlineRefTarget{sga3:iii:remark:3-6}\textbf{Remark 3.6.}
Let \(X_{\mathscr J}\) be a smooth affine \(S_{\mathscr J}\)-scheme.
By \SGARefLink{sga3:iii:proposition:0-15}{0.15}, up to isomorphism there is a unique \(S\)-scheme \(X\), smooth
over \(S\), whose reduction modulo \(\mathscr J\) is
\(X_{\mathscr J}\). If \(X_{\mathscr J}\) is endowed with an
\(S_{\mathscr J}\)-group structure, condition \((i_1)\) is automatic
by \SGARefLink{sga3:iii:proposition:0-16}{0.16}. Moreover, \SGARefLink{sga3:iii:remark:0-6}{0.6} simplifies the definition of \(L_0\) and gives:

\medskip
\noindent\SGAInlineRefTarget{sga3:iii:corollary:3-7}\textbf{Corollary 3.7.}
\textit{Let \(S,\mathscr I,\mathscr J\) be as in \SGARefLink{sga3:iii:lemma:3-1}{3.1}, and let
\(X_{\mathscr J}\) be a smooth affine \(S_{\mathscr J}\)-group.}
\begin{enumerate}
\item[(i)] \textit{The set of smooth \(S\)-groups whose reduction
modulo \(\mathscr J\) is \(X_{\mathscr J}\), up to isomorphism
inducing the identity on \(X_{\mathscr J}\), is either empty or
principal homogeneous under}
\[
 H^2\!\left(
  X_0,\operatorname{Lie}(X_0/S_0)
  \otimes_{\mathscr O_{S_0}}\mathscr J
 \right).
\]
\item[(ii)] \textit{A smooth \(S\)-group whose reduction modulo
\(\mathscr J\) is \(X_{\mathscr J}\) exists if and only if a certain
obstruction in}
\[
 H^3\!\left(
  X_0,\operatorname{Lie}(X_0/S_0)
  \otimes_{\mathscr O_{S_0}}\mathscr J
 \right)
\]
\textit{is zero.}
\end{enumerate}

As usual, one obtains the following corollaries.

\medskip
\noindent\SGAInlineRefTarget{sga3:iii:corollary:3-8}\textbf{Corollary 3.8.}
\textit{Let \(S\) be a scheme and \(S_0\) a closed subscheme defined by
a nilpotent ideal \(\mathscr I\). Let \(X_0\) be a smooth affine
\(S_0\)-group.}
\begin{enumerate}
\item[(i)] \textit{If, for every \(n\geq0\),}
\[
 H^2\!\left(
  X_0,\operatorname{Lie}(X_0/S_0)
  \otimes_{\mathscr O_{S_0}}
  \mathscr I^{n+1}/\mathscr I^{n+2}
 \right)=0,
\]
\textit{then any two smooth \(S\)-groups whose reduction is \(X_0\)
are isomorphic by an isomorphism inducing the identity on \(X_0\).}
\item[(ii)] \textit{If, for every \(n\geq0\),}
\[
 H^3\!\left(
  X_0,\operatorname{Lie}(X_0/S_0)
  \otimes_{\mathscr O_{S_0}}
  \mathscr I^{n+1}/\mathscr I^{n+2}
 \right)=0,
\]
\textit{then there exists a smooth \(S\)-group whose reduction is
\(X_0\).}
\end{enumerate}

\medskip
\noindent\SGAInlineRefTarget{sga3:iii:corollary:3-9}\textbf{Corollary 3.9.}
\textit{Let \(S\) be affine and \(S_0\) a closed subscheme defined by a
nilpotent ideal \(\mathscr I\). Suppose that every
\(\mathscr I^{n+1}/\mathscr I^{n+2}\) is locally free on \(S_0\). Let
\(X_0\) be a smooth affine \(S_0\)-group.}
\begin{enumerate}
\item[(i)] \textit{If
\(H^2(X_0,\operatorname{Lie}(X_0/S_0))=0\), then any two smooth
\(S\)-groups whose reduction is \(X_0\) are isomorphic.}
\item[(ii)] \textit{If
\(H^3(X_0,\operatorname{Lie}(X_0/S_0))=0\), then there exists a smooth
\(S\)-group whose reduction is \(X_0\).}
\end{enumerate}

\medskip
\noindent\SGAInlineRefTarget{sga3:iii:corollary:3-10}\textbf{Corollary 3.10.}
\textit{Let \(S_0\) be a scheme and let \(S=I_{S_0}\) be the
dual-number scheme over \(S_0\). Let \(X_0\) be a smooth \(S_0\)-group.
Then every smooth \(S\)-group \(Y\) for which \(Y_0\) is
\(S_0\)-isomorphic to \(X_0\) is \(S\)-isomorphic to
\(X=X_0\times_{S_0}S\) if and only if}
\[
 H^2(X_0,\operatorname{Lie}(X_0/S_0))=0.
\]
\footnote{Editor's note: This result is used in XXIV, \SGARefLink{sga3:xxiv:1.13}{1.13}.}

Indeed, by \SGARefLink{sga3:iii:theorem:3-5}{3.5} the set of classes of such groups \(Y\), up to an
\(S\)-group isomorphism inducing the identity on \(X_0\), is in
bijection with \(H^2(X_0,\operatorname{Lie}(X_0/S_0))\). Hence their
classes up to an arbitrary \(S\)-group isomorphism are in bijection
with
\[
 H^2(X_0,\operatorname{Lie}(X_0/S_0))/\Gamma_0,
\]
where
\[
 \Gamma_0=\operatorname{Aut}_{S_0\text{-}\mathrm{gr.}}(X_0),
\]
which acts in the evident way on \(H^2\). The conclusion follows at
once.
\footnote{Editor's note: Indeed,
\(\operatorname{Aut}_{S_0\text{-}\mathrm{gr.}}(X_0)\) acts by group
automorphisms on the abelian group
\(H^2(X_0,\operatorname{Lie}(X_0/S_0))\), so the orbit of \(0\) is the
singleton \(\{0\}\). Consequently, the quotient set is a singleton if
and only if
\(H^2(X_0,\operatorname{Lie}(X_0/S_0))=\{0\}\).}

